\subsubsection{Subproceso 2: Selección y Registro de Cursos}



En correspondencia con las técnicas empleadas, se utilizaron los siguientes instrumentos:

\textbf{Ficha de revisión documental:} Formato estructurado en el que se registraron, para cada fuente normativa consultada, el nombre del documento, su código o número de resolución, la fecha de emisión, los artículos o secciones relevantes y el extracto de las reglas de negocio identificadas (límites de créditos, requisitos mínimos, reglas de prerrequisito y de incompatibilidad horaria). Este instrumento permitió organizar y contrastar la información proveniente de múltiples documentos institucionales de manera sistemática.

\textbf{Guía de observación del sistema (lista de cotejo de pantallas):} Instrumento compuesto por una lista de pasos del flujo de selección de cursos en la Intranet DIRCE, con campos para registrar: el nombre de cada pantalla o módulo visitado, las acciones disponibles para el usuario, los mensajes de validación emitidos por el sistema, los campos visibles y los comportamientos observados (bloqueos, alertas, tiempos de carga). Complementariamente, se realizaron capturas de pantalla de las interfaces principales del portal para respaldar lo registrado en la guía.


